\section{The {\v C}ech complex and {\v C}ech cohomology}
\label{section-cech}

\noindent
Let $X$ be a topological space.
Let $\mathcal{U} : U = \bigcup_{i \in I} U_i$ be an open covering,
see Topology, Basic notion (\ref{topology-item-covering}).
As is customary we denote
$U_{i_0\ldots i_p} = U_{i_0} \cap \ldots \cap U_{i_p}$ for the
$(p + 1)$-fold intersection of members of $\mathcal{U}$.
Let $\mathcal{F}$ be an abelian presheaf on $X$.
Set
$$
\check{\mathcal{C}}^p(\mathcal{U}, \mathcal{F})
=
\prod\nolimits_{(i_0, \ldots, i_p) \in I^{p + 1}}
\mathcal{F}(U_{i_0\ldots i_p}).
$$
This is an abelian group. For
$s \in \check{\mathcal{C}}^p(\mathcal{U}, \mathcal{F})$ we denote
$s_{i_0\ldots i_p}$ its value in $\mathcal{F}(U_{i_0\ldots i_p})$.
Note that if $s \in \check{\mathcal{C}}^1(\mathcal{U}, \mathcal{F})$
and $i, j \in I$ then $s_{ij}$ and $s_{ji}$ are both elements
of $\mathcal{F}(U_i \cap U_j)$ but there is no imposed
relation between $s_{ij}$ and $s_{ji}$. In other words, we are {\it not}
working with alternating cochains (these will be defined
in Section \ref{section-alternating-cech}). We define
$$
d : \check{\mathcal{C}}^p(\mathcal{U}, \mathcal{F})
\longrightarrow
\check{\mathcal{C}}^{p + 1}(\mathcal{U}, \mathcal{F})
$$
by the formula
\begin{equation}
\label{equation-d-cech}
d(s)_{i_0\ldots i_{p + 1}}
=
\sum\nolimits_{j = 0}^{p + 1}
(-1)^j
s_{i_0\ldots \hat i_j \ldots i_{p + 1}}|_{U_{i_0\ldots i_{p + 1}}}
\end{equation}
It is straightforward to see that $d \circ d = 0$. In other words
$\check{\mathcal{C}}^\bullet(\mathcal{U}, \mathcal{F})$ is a complex.

\begin{definition}
\label{definition-cech-complex}
Let $X$ be a topological space.
Let $\mathcal{U} : U = \bigcup_{i \in I} U_i$ be an open covering.
Let $\mathcal{F}$ be an abelian presheaf on $X$.
The complex $\check{\mathcal{C}}^\bullet(\mathcal{U}, \mathcal{F})$
is the {\it {\v C}ech complex} associated to $\mathcal{F}$ and the
open covering $\mathcal{U}$. Its cohomology groups
$H^i(\check{\mathcal{C}}^\bullet(\mathcal{U}, \mathcal{F}))$ are
called the {\it {\v C}ech cohomology groups} associated to
$\mathcal{F}$ and the covering $\mathcal{U}$.
They are denoted $\check H^i(\mathcal{U}, \mathcal{F})$.
\end{definition}

\begin{lemma}
\label{lemma-cech-h0}
Let $X$ be a topological space.
Let $\mathcal{F}$ be an abelian presheaf on $X$.
The following are equivalent
\begin{enumerate}
\item $\mathcal{F}$ is an abelian sheaf and
\item for every open covering $\mathcal{U} : U = \bigcup_{i \in I} U_i$
the natural map
$$
\mathcal{F}(U) \to \check{H}^0(\mathcal{U}, \mathcal{F})
$$
is bijective.
\end{enumerate}
\end{lemma}

\begin{proof}
This is true since the sheaf condition is exactly that
$\mathcal{F}(U) \to \check{H}^0(\mathcal{U}, \mathcal{F})$
is bijective for every open covering.
\end{proof}

\begin{lemma}
\label{lemma-cech-trivial}
Let $X$ be a topological space. Let $\mathcal{F}$ be an abelian presheaf on $X$.
Let $\mathcal{U} : U = \bigcup_{i \in I} U_i$ be an open covering. If
$U_i = U$ for some $i \in I$, then the extended {\v C}ech complex
$$
\mathcal{F}(U) \to \check{\mathcal{C}}^\bullet(\mathcal{U}, \mathcal{F})
$$
obtained by putting $\mathcal{F}(U)$ in degree $-1$ with differential given by
the canonical map of $\mathcal{F}(U)$ into
$\check{\mathcal{C}}^0(\mathcal{U}, \mathcal{F})$
is homotopy equivalent to $0$.
\end{lemma}

\begin{proof}
Fix an element $i \in I$ with $U = U_i$. Observe that
$U_{i_0 \ldots i_p} = U_{i_0 \ldots \hat i_j \ldots i_p}$ if $i_j = i$.
Let us define a homotopy
$$
h :
\prod\nolimits_{i_0 \ldots i_{p + 1}} \mathcal{F}(U_{i_0 \ldots i_{p + 1}})
\longrightarrow
\prod\nolimits_{i_0 \ldots i_p} \mathcal{F}(U_{i_0 \ldots i_p})
$$
by the rule
$$
h(s)_{i_0 \ldots i_p} = s_{i i_0 \ldots i_p}
$$
In other words, $h : \prod_{i_0} \mathcal{F}(U_{i_0}) \to \mathcal{F}(U)$
is projection onto the factor $\mathcal{F}(U_i) = \mathcal{F}(U)$
and in general the map $h$ equals the projection onto the factors
$\mathcal{F}(U_{i i_1 \ldots i_{p + 1}}) =
\mathcal{F}(U_{i_1 \ldots i_{p + 1}})$.
We compute
\begin{align*}
(dh + hd)(s)_{i_0 \ldots i_p}
& =
\sum\nolimits_{j = 0}^p
(-1)^j
h(s)_{i_0 \ldots \hat i_j \ldots i_p}
+
d(s)_{i i_0 \ldots i_p}\\
& =
\sum\nolimits_{j = 0}^p
(-1)^j
s_{i i_0 \ldots \hat i_j \ldots i_p}
+
s_{i_0 \ldots i_p}
+
\sum\nolimits_{j = 0}^p
(-1)^{j + 1}
s_{i i_0 \ldots \hat i_j \ldots i_p} \\
& =
s_{i_0 \ldots i_p}
\end{align*}
This proves the identity map is homotopic to zero as desired.
\end{proof}






